%% kapa: exam=1 round=25 year=2014 period=2 form=A subject=accounting subject_ko=회계학 seq=40 number=80
%% source: Q-Net 감정평가사 기출문제 (cst00309) artlSeq=5101207 file=extracted/1008/5101207/2014년 제25회 감정평가사 1차시험 기출문제/2교시(A형).hwp
%% conversion: kordoc 4.24.51 -> tools/gen1.py (auto); review=proofread-sonnet; flags=none
\begin{question}{80}
(주)대한은 완제품 생산에 필요한 A부품을 매월 500단위씩 자가제조하고 있다. 그런데 타 회사에서 매월 A부품 500단위를 단위당 ￦100에 납품하겠다고 제의하였다. A부품을 자가제조할 경우 변동제조원가는 단위당 ￦70이고, 월간 고정제조간접원가 총액은 ￦50,000이다. 만약 A부품을 외부구입하면 변동제조원가는 발생하지 않으며, 월간 고정제조간접원가의 40\%를 절감할 수 있다. 또한 A부품 생산에 사용되었던 설비는 여유설비가 되며 다른 회사에 임대할 수 있다. A부품을 외부 구입함으로써 매월 ￦10,000의 이익을 얻고자 한다면, 여유설비의 월 임대료를 얼마로 책정해야 하는가?
\begin{choices}
  \item ￦5,000
  \item ￦6,000
  \item ￦7,000
  \item ￦8,000
  \item ￦10,000
\end{choices}
\end{question}
\answer{1}
